%-*-latex-*-
%-*-latex-*-
\newcommand\COURSE{ciss245}
\newcommand\ASSESSMENT{q03}
\newcommand\ASSESSMENTTYPE{Quiz}
\newcommand\POINTS{\textwhite{xxx/xxx}}

\input{myquizpreamble}
\input{yliow}
\input{\COURSE}
%-*-latex-*-

%-*-latex-*-

%-*-latex-*-
\renewcommand\TITLE{\ASSESSMENTTYPE \ \ASSESSMENT}

\renewcommand\EMAIL{}
\renewcommand\AUTHOR{eagraham1@cougars.ccis.edu}

\textwidth=6in
\begin{document}
\topmatterthree

%------------------------------------------------------------------------------
\nextq
Write a function that computes the sum of the
values in array \verb!x! from index 0 up to
and including index \verb!size - 1!.
The function prototype is
\begin{console}
double sum(double x[], int size);
\end{console}
\ANSWER
\begin{answercode}
double sum(double x[], int size)
{
  double num = 0;
  for(int i = 1; i < size; i++)
  {
    num += x[i];
  }
  return num;
}
\end{answercode}

%------------------------------------------------------------------------------
\nextq
Write a function that
computes the maximum of the
values in array \verb!x! from index 0 up to
and including index \verb!size - 1!.
The function prototype is
\begin{console}
int max(int x[], int size);
\end{console}
\ANSWER
\begin{answercode}
int max(int x[], int size)
{
  int max = 0;
  for(int i = 0; i < size; i++)
  {
    if(max < x[i])
    {
      max = x[i];
    }
  }
  return max;
}
\end{answercode}

%------------------------------------------------------------------------------
\nextq
Write a function that performs a linear search
for \verb!target! in array \verb!x! from index 0 up to
and including index \verb!size - 1!.
If \verb!target! is not found, \verb!-1! is returned.
The function prototype is
\begin{console}
int linearsearch(int x[], int size, int target);
\end{console}
\ANSWER
\begin{answercode}
int linearseach(int x[], int size, int target)
{
  for(int i = 0; i < size; i++)
  {
    if(x[i] == target)
    {
      return i;
    }
  }
  return -1;
}
\end{answercode}

%------------------------------------------------------------------------------
\nextq
Write a function that counts the number of times
\verb!target! appears in
array \verb!x! from index 0 up to
and including index \verb!size - 1!.
The function prototype is
\begin{console}
int count(int x[], int size, int target);
\end{console}
\ANSWER
\begin{answercode}
int count(int x[], int size, int target)
{
  int c = 0;
  for(int i = 0; i < size; i++)
  {
    if(x[i] == target)
    {
      c++;
    }
  }
  return c;
}
\end{answercode}

%------------------------------------------------------------------------------
\nextq
Write a function that swaps the values of two integer variables.
Part of the function prototype is given:
\begin{console}
    swap(     x,    y);
\end{console}
\ANSWER
\begin{answercode}
void swap(int & x, int & y)
{
  int temp = y;
  y = x;
  x = temp;
}
\end{answercode}

%------------------------------------------------------------------------------
%-*-latex-*-
\newpage
\input{instructions}
\end{document}

